\begin{sloppypar*}
    The proposed solution\cite{deepSimDEFOFCLimpl} sets up a supervised neural 
    network which takes in GO-term embeddings (\textit{from some other source}),
    utilizes all sub-ontologies (BP, MF, CC) either individually (\textbf{single
    channel}) or simultaneously (\textbf{multi channel}) and performs \textit{classification}
    or \textit{regression} for some downstream biological task, generating learned
    embeddings in the due course. The main components of the logic presented in 
    the paper may be visualized through the \textit{mindmap} in Figure \ref{fig:mindmap};
    wherein, under \textbf{Training Model}, \verb|SCN| means \textit{Single Channel Network}
    and \verb|MCN| means \textit{Multi-Channel Network}.\hfill\newline

    % \begin{markdown}
    % \end{markdown}
    \begin{figure}
        \centering
        \includegraphics[width=170mm,scale=1]{mindmap.png}
        \caption{Mindmap}
        \label{fig:mindmap}
    \end{figure}

    Three downstream tasks have been used, where, \textbf{Protein-Protein Interactions}
    (PPI) has been setup as a \textit{classification} task and \textbf{Gene Expression}
    (GE) and \textbf{Sequence Homology} (SH) have been setup as \textit{regression} tasks.

    \subsection{Experimental Data}
        The data for the experimental setup has been acquired from the sources as described below:
        \begin{enumerate}
            \item The \textbf{Gene Ontology} information has been downloaded from 
            \verb|www.geneontology.org/page/download-ontology|
            \item The term definition texts were enriched using the \textbf{MEDLINE}
            abstracts downloaded from \verb|www.nlm.nih.gov/databases/download/data_distrib_main.html|
            \item For PPI classification downstream task, \textbf{STRING} dataset
            was obtained from \verb|string-db.org/cgi/download.pl|
            \item For SH regression downstream task, \textbf{UniProt in FASTA format}
            was obtained from \verb|www.uniprot.org/proteomes/*|, where, * is replaced by a 
            \textit{code} for either \textbf{\textit{yeast}} or for \textbf{\textit{human}}
            \item For GE regression downstream task, for \textit{yeast}, data was
            obtained from \textit{supplementary material} of \cite{geneexpryeast} and for
            \textit{human}, \textbf{GTEx} data was obtained from \verb|www.gtexportal.org|
        \end{enumerate}
    
    \subsection{Pre-training Embeddings}
        \begin{figure}
            \centering
            \includegraphics[scale=0.3]{pretrain_embed.png}
            \caption{Pre-training Embedding}
            \label{fig:pretrain}
        \end{figure}
        In Figure \ref{fig:pretrain}, the overall process for creating the input vectors 
        for the neural network has been depicted via a \textit{flowchart}. In the figure,
        the term \verb|SOC| means \textit{Second Order Co-occurrence} (matrix), \verb|PMI|
        means \textit{Pointwise Mutual Information}, \verb|LSA| means \textit{Latent Semantic Analysis}
        and \verb|SVD| means \textit{Singular Value Decomposition}.\hfill\newline

        \noindent The purpose of the initialization of the embeddings is to ensure that
        gene products have good initial representation in the neural network so that meaningful
        gradients are generated.

    \subsection{Training Model}
        Whereas, the layers in the \textit{Single Channel} and \textit{Multi Channel}
        versions have been summarised in the \textit{mindmap} in Figure \ref{fig:mindmap},
        in this section we will capture the logic of the neural network via the \textit{flowchart}
        in Figure \ref{fig:network}.

        \begin{figure}
            \centering
            \includegraphics[scale=0.243]{training_network_mc.png}
            \caption{Training Model}
            \label{fig:network}
        \end{figure}

        \noindent In the Figure \ref{fig:network} the data sources indicate the
        different downstream tasks that have already been discussed above. The process
        may be further enumerated with certain more details as follows (\textit{the multi-channel
        process is described}):

        \begin{enumerate}
            \item A data source for a particular task is selected, and pairs of 
            relevant gene-products are selected
            \item Each gene product may contain multiple GO terms. Further, the 
            gene terms are different under each sub-ontology (BP, CC and MF). So, 3 sets
            of GO terms are formed for each gene-product
            \item The term embeddings are looked-up in the tables formed in the
            pre-training embedding step, and matrices are formed. When the number of
            terms are not same, padding of large negative numbers are used
            \item The first maxpool layer selects the max value from the fixed 100 dimension
            features for each term and prepares a 1-D vector from the matrix
        \end{enumerate}

        \noindent The remaining steps are as per the \textit{flowchart}. In case of 
        the \textbf{Single Channel} version, the channel merge layer is not required.
        The \textbf{Gene Embedding} can be extracted right before the \textit{gene-pair
        merge layer}.
        
\end{sloppypar*}